\documentclass[12pt,a4paper]{article}
\usepackage[utf8]{inputenc}
\usepackage[T1]{fontenc}
\usepackage{geometry}
\geometry{margin=2.5cm}
\usepackage{amsmath}
\usepackage{booktabs}
\usepackage{hyperref}
\usepackage{graphicx}
\usepackage{array}
\usepackage{multirow}
\usepackage{setspace}
\usepackage{parskip}
\usepackage{titlesec}
\usepackage{natbib}
\usepackage{xcolor}

\titleformat{\section}{\large\bfseries}{\thesection}{1em}{}
\titleformat{\subsection}{\normalsize\bfseries}{\thesubsection}{1em}{}

\hypersetup{
    colorlinks=true,
    linkcolor=blue,
    citecolor=blue,
    urlcolor=blue
}

\title{\textbf{Mining Behavioral Risk Patterns for Diabetes Prediction:\\An Explainable Ensemble Approach}}
\author{
    Group 9\\
    \small Class: CLC01\\
    \small Course: Big Data Mining\\
    \small Hanoi University --- Faculty of Information Technology
}
\date{April 2026}

\begin{document}
\maketitle
\tableofcontents
\newpage

%% ============================================================
\begin{abstract}
This paper presents a hybrid ARM--SHAP ensemble framework for type 2 diabetes mellitus
(T2DM) prediction applied to the BRFSS 2015 dataset (269,131 records). The study
synthesizes five methodological pillars: association rule mining (ARM), tree-based ensemble
classifiers, class-imbalance handling, SHAP-based explainable AI (XAI), and epidemiological
context for Vietnam. The central finding is a persistent methodological gap in existing
literature: studies excel in either unsupervised pattern discovery through ARM or post-hoc
model interpretability through SHAP, but none implement a sequential pipeline that generates
behavioral risk hypotheses via FP-Growth, trains a hyperparameter-tuned ensemble classifier,
and verifies genuine interaction effects through targeted SHAP subgroup analysis on the same
dataset. Our pipeline addresses this gap: FP-Growth yielded 157 high-confidence rules
(lift $\geq$ 1.5); Optuna-tuned XGBoost achieved AUC-ROC = 0.848 and recall = 0.814;
SHAP subgroup verification confirmed 2 of 4 ARM-derived hypotheses; and age-stratified
analysis revealed that mental health burden is 1.69$\times$ more predictive in adults aged
18--44 than in the general population, with implications for T2DM prevention in Vietnam.
\end{abstract}

\textbf{Keywords:} association rule mining; diabetes prediction; ensemble learning;
explainable AI; SHAP explainability; BRFSS

%% ============================================================
\section{Introduction}

\subsection{Background and Motivation}

Type 2 diabetes mellitus (T2DM) continues to escalate as a major non-communicable disease
worldwide. The International Diabetes Federation estimates 589 million adults living with
diabetes globally in 2024, with projections reaching approximately 853 million by 2050;
over 90\% of cases are type 2 and largely driven by modifiable behavioral factors
\citep{Duncan2025}. In Vietnam, the disease affects an estimated 2.5 million adults,
with age-standardized prevalence around 3.4\% and approximately 37.8\% of cases remaining
undiagnosed, with incidence rising most rapidly among adults under 45 amid urbanization
and lifestyle changes \citep{Vuong2024, Nguyen2015}. Traditional epidemiological approaches
identify individual risk factors but rarely capture synergistic multi-feature behavioral
patterns that drive disease progression.

Two methodological traditions dominate recent machine-learning research on T2DM prediction.
The first combines ensemble tree-based models with SHAP for interpretability, yielding both
accurate predictions and post-hoc feature attributions \citep{Majcherek2025, Kutlu2024,
Netayawijit2025, Ahmed2024, Rafie2025, Islam2024}. The second employs association rule
mining to discover co-occurrence patterns among behavioral variables \citep{Fakir2024,
Bata2025}. These streams have largely remained separate: ARM studies generate multi-feature
hypotheses but cannot verify whether co-occurrences represent genuine risk amplification,
while SHAP-ensemble studies provide rigorous feature explanations but lack prior hypotheses
about which combinations to examine.

\subsection{Problem Statement}

Existing T2DM prediction studies suffer from five critical gaps. First, ARM studies generate
co-occurrence patterns without mechanisms to verify genuine risk-amplifying interactions.
Second, SHAP studies apply explainability as a post-hoc tool without prior testable
hypotheses. Third, feature leakage from variables such as \texttt{GenHlth} (r = 0.33),
\texttt{CholCheck}, and \texttt{DiffWalk} remains unaddressed in BRFSS-based studies.
Fourth, no study performs age-stratified SHAP analysis despite evidence that T2DM rises
fastest among younger adults. Fifth, no study implements the full sequential pipeline of
ARM hypothesis generation, Optuna-tuned ensemble prediction, and targeted SHAP subgroup
verification on the same large-scale behavioral dataset.

\subsection{Research Questions}

This study addresses three research questions:
\begin{enumerate}
    \item Which multi-feature combinations of behavioral and biological indicators constitute
    the strongest risk patterns for T2DM?
    \item Do the behavioral risk combinations identified through ARM genuinely amplify
    individual feature contributions, or do they merely reflect statistical co-occurrence?
    \item Do the behavioral features driving risk differ between the younger population
    (ages 18--44) and the general adult population?
\end{enumerate}

\subsection{Contributions}

The main contributions of this work are: (1) a sequential ARM-to-SHAP pipeline that
generates testable hypotheses before applying explainability analysis; (2) systematic
feature leakage quantification through three model variants; (3) Optuna-tuned ensemble
classification achieving AUC-ROC = 0.848 on BRFSS 2015; (4) SHAP subgroup verification
confirming that 2 of 4 ARM-derived hypotheses represent genuine interaction effects;
and (5) age-stratified SHAP analysis revealing disproportionate mental health and income
risk amplification in adults aged 18--44, with implications for Vietnam.

%% ============================================================
\section{Related Work}

\subsection{Association Rule Mining in Healthcare}

Association rule mining (ARM) is a core unsupervised technique for uncovering hidden
relationships in categorical data. \citet{Agrawal1994} introduced the Apriori algorithm,
which generates candidate itemsets level-wise and prunes infrequent sets using a minimum-support
threshold. While foundational, Apriori suffers from repeated full-database scans and explosive
candidate growth in high-dimensional spaces. \citet{Han2000} addressed these limitations via
FP-Growth, which constructs a compact FP-tree to mine frequent patterns without candidate
generation, making it well-suited for large transactional datasets such as BRFSS.

In healthcare contexts, \citet{Bata2025} employed Apriori-based ARM on time-stratified
diabetes comorbidity data, confirming ARM's value for hypothesis generation but noting that
rule quality degrades on imbalanced data. \citet{Fakir2024} implemented parallelized FP-Growth
on Apache Spark, linking high-calorie intake, low physical activity, and elevated BMI with
diabetes onset. Critically, neither study verified whether discovered patterns represent
genuine risk-amplifying interactions.

\subsection{Ensemble Classifiers and Hyperparameter Optimization}

Tree ensembles dominate tabular health-data tasks. \citet{Breiman2001} introduced Random
Forests, reducing variance through bagging and random feature selection. \citet{Chen2016}
developed XGBoost with L1/L2 regularization, consistently ranking among top performers on
imbalanced health benchmarks. \citet{Ke2017} proposed LightGBM with Gradient-based One-Side
Sampling (GOSS), achieving substantial speed improvements through leaf-wise growth.
Hyperparameter tuning via Optuna \citep{Akiba2019}, a Bayesian optimization framework using
Tree-structured Parzen Estimators, enables systematic exploration of high-dimensional
hyperparameter spaces while providing reproducibility.

\subsection{SHAP-Based Explainable AI for Diabetes Prediction}

\citet{Lundberg2017} introduced SHAP (SHapley Additive exPlanations), a game-theoretic
approach assigning consistent feature attributions. \citet{Lundberg2019} extended this to
Tree SHAP for exact polynomial-time computation on tree ensembles. \citet{Ahmed2024}
empirically confirmed SHAP's superiority over LIME for tree-based models.

Two studies are directly comparable to the current work. \citet{Majcherek2025} benchmarked
18 models on BRFSS 2015, identifying HighBP, BMI, Age, and GenHlth as dominant contributors,
but applied SHAP post-hoc without prior hypotheses and retained GenHlth (r = 0.33) without
leakage analysis. \citet{Kutlu2024} applied XGBoost with Recursive Feature Elimination,
achieving 86.6\% accuracy with SHAP-confirmed BMI thresholds around 30 kg/m$^2$, but also
lacked systematic leakage analysis. Neither study performs age-stratified analysis.

\subsection{Class Imbalance Handling}

\citet{Chawla2002} introduced SMOTE, generating synthetic minority examples by interpolating
between nearest neighbors. \citet{Agyemang2025} demonstrated 8--15\% sensitivity improvements
from oversampling variants. A critical design consideration is SMOTE-aware cross-validation:
SMOTE must be fitted only on training folds to prevent information leakage from validation
instances \citep{Netayawijit2025, Rezki2024}.

\subsection{Epidemiological Context for Vietnam}

\citet{Duncan2025} documented 589 million adults with diabetes globally in 2024.
\citet{Vuong2024} reported age-standardized prevalence of approximately 3.4\% in Vietnam,
with 37.8\% of cases undiagnosed. \citet{Nguyen2015} documented that T2DM incidence is
rising fastest among Vietnamese adults under 45, driven by urbanization and dietary
transitions---directly motivating the age-stratified analysis in this study.

%% ============================================================
\section{Methodology}

\subsection{Dataset}

The study uses the BRFSS 2015 dataset (Behavioral Risk Factor Surveillance System),
a U.S. CDC telephone survey comprising 269,131 records and 22 features covering
behavioral, clinical, and demographic variables. After binarization (merging pre-diabetes
and diabetes into At-Risk), the class distribution is 72.22\% No-Risk and 27.78\% At-Risk.

\subsection{Data Preparation}

The preparation pipeline includes: (1) target binarization; (2) BMI outlier handling via
IQR-capping at 42.5 (affecting 11,280 records, vs.\ removal of BMI $>$ 60 affecting only
993 records---IQR-capping was selected to preserve data); (3) MinMaxScaler normalization
fitted on training data only; and (4) discretization of continuous/ordinal features into
categorical bins for ARM (\texttt{BMI\_cat}, \texttt{Age\_cat}, \texttt{GenHlth\_cat},
\texttt{MentHlth\_cat}, \texttt{PhysHlth\_cat}, \texttt{Income\_cat},
\texttt{Education\_cat}).

\subsection{Association Rule Mining via FP-Growth}

FP-Growth (mlxtend) is applied to the transaction matrix (41 binary columns after one-hot
encoding) with minimum support = 0.05, minimum confidence = 0.45, and minimum lift = 1.5.
Rules with consequent = AtRisk are retained and ranked by lift. The top rules are used to
formulate testable hypotheses about SHAP contributions within ARM-defined subgroups:

\begin{equation}
H_i: \bar{\phi}_f^{(\text{sub})} > \bar{\phi}_f^{(\text{pop})}
\end{equation}

where $\bar{\phi}_f^{(\text{sub})}$ is the mean SHAP value of feature $f$ within the
ARM-defined subgroup, and $\bar{\phi}_f^{(\text{pop})}$ is the population mean SHAP value.
A hypothesis is confirmed if the relative increase exceeds 50\%.

\subsection{Classification Pipeline}

Six classifiers are trained using 5-fold Stratified K-Fold cross-validation:
Logistic Regression and Decision Tree (baselines), Random Forest, XGBoost, and LightGBM
(each tuned via Optuna with 50 trials optimizing macro-F1), and a Soft Voting Ensemble
adopted only if it outperforms the best single model on macro-F1.

Class imbalance is addressed through three strategies: SMOTE-aware cross-validation
(SMOTE fitted only on training folds using imbalanced-learn Pipeline), NearMiss, and
class\_weight balancing. The best strategy is selected by mean macro-F1 on validation folds.

\subsection{Feature Leakage Analysis}

Three model variants quantify leakage sensitivity:
\begin{itemize}
    \item \textbf{Variant A} (full 21 features): includes all features
    \item \textbf{Variant B} (leakage-removed): excludes GenHlth, CholCheck, DiffWalk,
    HeartDiseaseorAttack, Stroke
    \item \textbf{Variant C} (behavioral-only): PhysActivity, BMI, Smoker, Fruits, Veggies,
    HvyAlcoholConsump, Income, Education, Age, Sex
\end{itemize}

\subsection{SHAP Explainability and Subgroup Verification}

SHAP TreeExplainer computes attributions on the best model (XGBoost Variant A).
Targeted subgroup analysis verifies each ARM-derived hypothesis by comparing mean SHAP
values within ARM-defined subgroups against the population mean. Youth-versus-general
SHAP comparison (Age codes 1--5, ages 18--44) addresses RQ3.

All experiments use Python 3.12 with scikit-learn, XGBoost, LightGBM, SHAP, Optuna,
and mlxtend (seed = 42).

%% ============================================================
\section{Implementation and Results}

\subsection{Association Rule Mining}

FP-Growth yielded 19,204 frequent itemsets and 10,346 candidate rules. After filtering
for consequent = AtRisk and confidence $\geq$ 0.45, 157 rules remained. The dataset-wide
At-Risk base rate is 27.78\%, so all reported rules represent meaningful lift above chance.

The highest-lift rule---\{CholCheck, DiffWalk, HighBP, HighChol\} $\rightarrow$ AtRisk---achieved
confidence = 0.610 and lift = 2.194. Among rules with at least one behavioral feature,
the strongest was \{BMI\_cat\_SeverelyObese, HighBP\} $\rightarrow$ AtRisk (support = 0.056,
confidence = 0.585, lift = 2.105). Four hypotheses were formulated for SHAP verification:

\begin{itemize}
    \item \textbf{H1}: In subgroup \{AnyHealthcare, CholCheck, HighBP\}, mean SHAP of BMI
    $>$ population mean SHAP (rule lift = 2.124)
    \item \textbf{H2}: In subgroup \{CholCheck, HighChol, HighBP\}, mean SHAP of DiffWalk
    $>$ population mean SHAP (rule lift = 2.194)
    \item \textbf{H3}: In subgroup \{AnyHealthcare, CholCheck, GenHlth\_Fair\}, mean SHAP
    of HighChol $>$ population mean SHAP (rule lift = 2.062)
    \item \textbf{H4}: In subgroup \{AnyHealthcare, CholCheck, HighBP\}, mean SHAP of
    GenHlth $>$ population mean SHAP (rule lift = 2.049)
\end{itemize}

\subsection{Model Performance}

SMOTE outperformed NearMiss (macro-F1 = 0.692 vs.\ 0.590) and class\_weight (macro-F1 = 0.692)
and was selected as the primary balancing strategy. Table~\ref{tab:models} presents full
results on the held-out test set (n = 53,827).

\begin{table}[h]
\centering
\caption{Model performance on BRFSS 2015 test set (n = 53,827)}
\label{tab:models}
\begin{tabular}{lcccc}
\toprule
\textbf{Model} & \textbf{Recall} & \textbf{Macro-F1} & \textbf{AUC-ROC} & \textbf{MCC} \\
\midrule
Logistic Regression       & 0.746 & 0.694 & 0.807 & 0.419 \\
Decision Tree             & 0.793 & 0.784 & 0.822 & 0.578 \\
XGBoost (Optuna-tuned)    & \textbf{0.814} & \textbf{0.730} & \textbf{0.848} & \textbf{0.496} \\
LightGBM (Optuna-tuned)   & 0.613 & 0.711 & 0.814 & 0.423 \\
Random Forest (Optuna-tuned) & 0.700 & 0.717 & 0.824 & 0.445 \\
Soft Voting Ensemble      & 0.733 & 0.734 & 0.844 & 0.480 \\
\bottomrule
\end{tabular}
\end{table}

XGBoost achieved the highest AUC-ROC (0.848) and recall (0.814) and was retained as the
primary model. The Soft Voting Ensemble did not surpass XGBoost on macro-F1 (0.734 vs.\
0.730), so XGBoost was used per the pre-specified selection criterion.

\subsection{Feature Leakage Sensitivity}

Table~\ref{tab:leakage} shows the impact of removing potentially leaky features.

\begin{table}[h]
\centering
\caption{Leakage sensitivity analysis (XGBoost)}
\label{tab:leakage}
\begin{tabular}{lccc}
\toprule
\textbf{Variant} & \textbf{Recall} & \textbf{Macro-F1} & \textbf{AUC-ROC} \\
\midrule
A -- Full (21 features)       & 0.814 & 0.730 & 0.848 \\
B -- Leakage-removed          & 0.807 & 0.704 & 0.825 \\
C -- Behavioral-only          & 0.838 & 0.612 & 0.759 \\
\bottomrule
\end{tabular}
\end{table}

Removing leakage-suspect features (Variant B) reduced AUC by 0.023 and macro-F1 by 0.026,
confirming a modest but non-trivial leakage effect. Variant C's substantially lower
macro-F1 (0.612) demonstrates that purely behavioral features are insufficient for
high-accuracy prediction.

\subsection{SHAP Feature Importance}

Table~\ref{tab:shap} presents the top-10 features by mean $|\text{SHAP}|$ value from
XGBoost Variant A.

\begin{table}[h]
\centering
\caption{Top 10 features by mean absolute SHAP value (XGBoost, Variant A)}
\label{tab:shap}
\begin{tabular}{clc}
\toprule
\textbf{Rank} & \textbf{Feature} & \textbf{Mean $|\text{SHAP}|$} \\
\midrule
1  & GenHlth   & 0.674 \\
2  & Age       & 0.619 \\
3  & BMI       & 0.616 \\
4  & HighBP    & 0.515 \\
5  & HighChol  & 0.345 \\
6  & Income    & 0.264 \\
7  & MentHlth  & 0.229 \\
8  & PhysHlth  & 0.206 \\
9  & Education & 0.171 \\
10 & Sex       & 0.170 \\
\bottomrule
\end{tabular}
\end{table}

GenHlth, Age, and BMI collectively dominate model predictions, consistent with
\citet{Majcherek2025} and \citet{Kutlu2024}. SHAP dependence plots revealed non-linear
BMI thresholds, with risk contributions increasing sharply above BMI $\approx$ 30 kg/m$^2$.

\subsection{ARM Hypothesis Verification}

Table~\ref{tab:hyp} summarizes the four hypothesis tests using a 50\% relative increase
threshold for confirmation.

\begin{table}[h]
\centering
\caption{SHAP subgroup verification of ARM-derived hypotheses}
\label{tab:hyp}
\begin{tabular}{llcccc}
\toprule
\textbf{H} & \textbf{Feature} & \textbf{Pop.\ Mean SHAP} & \textbf{Sub.\ Mean SHAP} & \textbf{Rel.\ Increase} & \textbf{Confirmed} \\
\midrule
H1 & BMI      & $-$0.402 & $-$0.261 & 35.1\% & \texttimes \\
H2 & DiffWalk & $-$0.013 & $-$0.004 & 67.8\% & \checkmark \\
H3 & HighChol & $-$0.073 & $-$0.064 & 13.1\% & \texttimes \\
H4 & GenHlth  & $-$0.368 & $-$0.098 & 73.4\% & \checkmark \\
\bottomrule
\end{tabular}
\end{table}

H2 and H4 were confirmed: DiffWalk's SHAP contribution increased by 67.8\% within the
CholCheck--HighChol--HighBP subgroup; GenHlth's contribution increased by 73.4\% within
the AnyHealthcare--CholCheck--HighBP subgroup. H1 and H3 were rejected (35.1\% and 13.1\%).

\subsection{Youth vs.\ General Population SHAP Comparison}

The youth subgroup (ages 18--44, n = 10,377, 9.1\% At-Risk) showed a markedly different
risk profile. Table~\ref{tab:youth} presents the comparison.

\begin{table}[h]
\centering
\caption{Mean $|\text{SHAP}|$ comparison: full population vs.\ youth subgroup}
\label{tab:youth}
\begin{tabular}{lccc}
\toprule
\textbf{Feature} & \textbf{Full Population} & \textbf{Youth (18--44)} & \textbf{Ratio} \\
\midrule
GenHlth   & 0.674 & 0.919 & 1.36$\times$ \\
Age       & 0.618 & 1.949 & 3.15$\times$ \\
BMI       & 0.616 & 0.747 & 1.21$\times$ \\
HighBP    & 0.515 & 0.570 & 1.11$\times$ \\
HighChol  & 0.345 & 0.435 & 1.26$\times$ \\
Income    & 0.264 & 0.375 & 1.42$\times$ \\
MentHlth  & 0.229 & 0.387 & 1.69$\times$ \\
PhysHlth  & 0.206 & 0.267 & 1.30$\times$ \\
Education & 0.171 & 0.229 & 1.34$\times$ \\
Sex       & 0.170 & 0.159 & 0.94$\times$ \\
\bottomrule
\end{tabular}
\end{table}

Age's SHAP contribution in the youth subgroup is 3.15$\times$ higher than in the full
population, and MentHlth showed a 1.69$\times$ amplification, indicating that mental
health burden is a disproportionately important risk factor for younger adults.

%% ============================================================
\section{Discussion and Conclusion}

\subsection{RQ1: Multi-Feature Behavioral Risk Patterns}

FP-Growth identified 157 high-confidence rules pointing to T2DM risk. The dominant patterns
cluster around comorbidity combinations---particularly DiffWalk, HighBP, HighChol, and
CholCheck---rather than purely behavioral features. Among behavioral rules, the
\{BMI\_cat\_SeverelyObese, HighBP\} combination (lift = 2.105) is the strongest purely
modifiable risk pattern, reinforcing the obesity--hypertension--diabetes triad
\citep{Duncan2025, Vuong2024}. The Jaccard similarity of 0.308 between ARM and SHAP
feature sets indicates moderate convergence, with four features in common: BMI, GenHlth,
HighBP, and HighChol.

\subsection{RQ2: Genuine Interaction Effects vs.\ Statistical Co-occurrence}

Two of four ARM-derived hypotheses were confirmed by SHAP subgroup analysis. H2's
confirmation (DiffWalk SHAP +67.8\% in the CholCheck--HighChol--HighBP subgroup) suggests
that difficulty walking genuinely amplifies risk prediction within a cluster of cardiovascular
and metabolic comorbidities. H4's confirmation (GenHlth SHAP +73.4\%) indicates that
self-rated fair health carries substantially greater predictive weight when co-occurring
with healthcare access and cholesterol monitoring.

The two rejected hypotheses are equally informative. H1's rejection (BMI SHAP increase
only 35.1\%) suggests that BMI and HighBP may capture overlapping variance rather than
genuinely interacting---consistent with the Tree SHAP correlation bias documented by
\citet{AlHudaibi2025}. These results validate the necessity of the ARM-to-SHAP verification
step, directly addressing the limitation identified in \citet{Bata2025} and \citet{Fakir2024}.

\subsection{Model Performance and Comparison}

XGBoost achieved AUC-ROC = 0.848 and recall = 0.814, broadly comparable to
\citet{Majcherek2025} and slightly below \citet{Kutlu2024}'s reported 86.6\% accuracy.
The current study's recall of 0.814 on the At-Risk class is substantially higher than
Kutlu et al.'s minority class recall, reflecting the benefit of SMOTE-aware cross-validation
\citep{Netayawijit2025, Rezki2024}. The leakage sensitivity analysis confirms that
GenHlth, CholCheck, and DiffWalk contribute approximately 0.023 AUC points, suggesting
prior BRFSS studies overestimated performance by a modest margin.

\subsection{RQ3: Age-Stratified Risk Drivers and Implications for Vietnam}

The youth subgroup analysis reveals a qualitatively different risk profile for adults aged
18--44. The 3.15$\times$ amplification of Age's SHAP contribution within this subgroup
suggests an accelerating risk trajectory in the late-30s to early-40s bracket, consistent
with \citet{Nguyen2015}. The 1.69$\times$ amplification of MentHlth is a novel finding:
mental health burden as a disproportionate risk factor for younger adults may reflect
stress-related metabolic dysregulation or depression-associated lifestyle deterioration.
Income's 1.42$\times$ amplification aligns with socioeconomic vulnerability being more
predictive in working-age populations \citep{Vuong2024}.

\subsection{Limitations}

Several limitations constrain generalizability. First, BRFSS is a U.S. telephone survey;
all findings should be treated as hypothesis-generating rather than directly prescriptive
for Vietnam. Second, the youth subgroup's low At-Risk prevalence (9.1\%, n = 946 cases)
limits statistical power. Third, the 50\% relative increase threshold for hypothesis
confirmation is operationally defined. Fourth, Tree SHAP's known correlation bias means
SHAP values for correlated features (BMI, HighBP, HighChol) should be interpreted with
caution \citep{Lundberg2019, AlHudaibi2025}.

\subsection{Conclusion}

This study presents a hybrid ARM--SHAP ensemble framework that bridges two previously
separate methodological traditions in T2DM prediction. The sequential pipeline---FP-Growth
hypothesis generation, Optuna-tuned XGBoost classification (AUC-ROC = 0.848), and targeted
SHAP subgroup verification---demonstrates that ARM rules do not uniformly represent genuine
risk-amplifying interactions: 2 of 4 hypotheses were confirmed, validating the necessity
of post-hoc verification. Age-stratified analysis reveals that mental health burden and
income are disproportionately predictive in adults aged 18--44, providing actionable
hypotheses for T2DM prevention programs targeting Vietnam's rapidly growing at-risk youth
demographic. Future work should validate these findings on Vietnamese clinical datasets
and explore causal inference methods to move beyond associative patterns.

%% ============================================================
\section*{Acknowledgments}
The BRFSS 2015 dataset is publicly available, collected and anonymized by the U.S.\ CDC.
No personally identifiable information is present. All findings are framed as statistical
associations, not causal claims. This research received no external funding. The authors
declare no conflicts of interest.

%% ============================================================
\bibliographystyle{plainnat}
\begin{thebibliography}{99}

\bibitem[Agrawal and Srikant(1994)]{Agrawal1994}
Agrawal, R., \& Srikant, R. (1994).
Fast algorithms for mining association rules.
\textit{Proceedings of the 20th VLDB Conference}, 487--499.

\bibitem[Agyemang et al.(2025)]{Agyemang2025}
Agyemang, E.~F., Mensah, J.~A., Nyarko, E., et al. (2025).
Addressing class imbalance problem in health data classification.
\textit{Applied Computational Intelligence and Soft Computing}, 2025(1).

\bibitem[Ahmed et al.(2024)]{Ahmed2024}
Ahmed, S., Kaiser, M.~S., Hossain, M.~S., \& Andersson, K. (2024).
A comparative analysis of LIME and SHAP interpreters with explainable ML-based diabetes predictions.
\textit{IEEE Access}, 13, 37370--37388.

\bibitem[Akiba et al.(2019)]{Akiba2019}
Akiba, T., Sano, S., Yanase, T., Ohta, T., \& Koyama, M. (2019).
Optuna: A next-generation hyperparameter optimization framework.
\textit{arXiv:1907.10902}.

\bibitem[Al-Hudaibi et al.(2025)]{AlHudaibi2025}
Al-Hudaibi, Z.~H., et al. (2025).
Explainable artificial intelligence-PREDICT.
\textit{Journal of Advanced Trends in Medical Research}, 2(2), 114--119.

\bibitem[Bata et al.(2025)]{Bata2025}
Bata, R., Ghanem, A.~S., Faludi, E.~V., Sztanek, F., \& Nagy, A.~C. (2025).
Association rule mining of time-based patterns in diabetes-related comorbidities on imbalanced data.
\textit{BMC Medical Informatics and Decision Making}, 25(1), 352.

\bibitem[Breiman(2001)]{Breiman2001}
Breiman, L. (2001).
Random forests.
\textit{Machine Learning}, 45(1), 5--32.

\bibitem[Chawla et al.(2002)]{Chawla2002}
Chawla, N.~V., Bowyer, K.~W., Hall, L.~O., \& Kegelmeyer, W.~P. (2002).
SMOTE: Synthetic minority over-sampling technique.
\textit{Journal of Artificial Intelligence Research}, 16, 321--357.

\bibitem[Chen and Guestrin(2016)]{Chen2016}
Chen, T., \& Guestrin, C. (2016).
XGBoost: A scalable tree boosting system.
\textit{arXiv:1603.02754}.

\bibitem[Duncan et al.(2025)]{Duncan2025}
Duncan, B.~B., Magliano, D.~J., \& Boyko, E.~J. (2025).
IDF Diabetes Atlas 11th edition 2025.
\textit{Nephrology Dialysis Transplantation}, 41(1), 7--9.

\bibitem[Fakir et al.(2024)]{Fakir2024}
Fakir, Y., Khalil, S., \& Fakir, M. (2024).
Extraction of association rules in a diabetic dataset using parallel FP-growth under Apache Spark.
\textit{International Journal of Informatics and Communication Technology}, 13(3), 445.

\bibitem[Han et al.(2000)]{Han2000}
Han, J., Pei, J., \& Yin, Y. (2000).
Mining frequent patterns without candidate generation.
\textit{ACM SIGMOD Record}, 29(2), 1--12.

\bibitem[Islam et al.(2024)]{Islam2024}
Islam, M.~M., et al. (2024).
Explainable machine learning for efficient diabetes prediction.
\textit{Engineering Reports}, 7(1).

\bibitem[Ke et al.(2017)]{Ke2017}
Ke, G., Meng, Q., Finley, T., et al. (2017).
LightGBM: A highly efficient gradient boosting decision tree.
\textit{Advances in Neural Information Processing Systems}, 30.

\bibitem[Kutlu et al.(2024)]{Kutlu2024}
Kutlu, M., Donmez, T.~B., \& Freeman, C. (2024).
Machine learning interpretability in diabetes risk assessment: A SHAP analysis.
\textit{Computers and Electronics in Medicine}, 1(1), 34--44.

\bibitem[Lundberg and Lee(2017)]{Lundberg2017}
Lundberg, S., \& Lee, S. (2017).
A unified approach to interpreting model predictions.
\textit{arXiv:1705.07874}.

\bibitem[Lundberg et al.(2019)]{Lundberg2019}
Lundberg, S.~M., Erion, G., Chen, H., et al. (2019).
Explainable AI for trees.
\textit{arXiv:1905.04610}.

\bibitem[Majcherek et al.(2025)]{Majcherek2025}
Majcherek, D., Ciesielski, A., \& Sobczak, P. (2025).
AI-driven analysis of diabetes risk determinants in U.S.\ adults.
\textit{PLoS ONE}, 20(9), e0328655.

\bibitem[Netayawijit et al.(2025)]{Netayawijit2025}
Netayawijit, P., Chansanam, W., \& Sorn-In, K. (2025).
Interpretable ML framework for diabetes prediction: Integrating SMOTE with SHAP.
\textit{Healthcare}, 13(20), 2588.

\bibitem[Nguyen et al.(2015)]{Nguyen2015}
Nguyen, C.~T., Pham, N.~M., Lee, A.~H., \& Binns, C.~W. (2015).
Prevalence of and risk factors for type 2 diabetes mellitus in Vietnam.
\textit{Asia Pacific Journal of Public Health}, 27(6), 588--600.

\bibitem[Rafie et al.(2025)]{Rafie2025}
Rafie, Z., et al. (2025).
Leveraging XGBoost and explainable AI for accurate prediction of type 2 diabetes.
\textit{BMC Public Health}, 25(1), 3688.

\bibitem[Rezki et al.(2024)]{Rezki2024}
Rezki, M.~K., et al. (2024).
Application of SMOTE to address class imbalance in diabetes classification.
\textit{Journal of Electronics Electromedical Engineering and Medical Informatics}, 6(4), 343--354.

\bibitem[Vuong et al.(2024)]{Vuong2024}
Vuong, T.~B., Tran, T.~M., \& Tran, N.~Q. (2024).
High prevalence of prediabetes and type 2 diabetes in high-risk adults in Vietnam.
\textit{Diabetes Epidemiology and Management}, 17, 100239.

\end{thebibliography}

\end{document}
